\bookchapter{บทที่ 24: เจาะลึกสิทธิ์ Root และกุญแจ BL}

\chapterintro{บทนี้เป็นการขยายต่อจากเนื้อหาบทที่ 20 อย่างลึกซึ้ง เหมาะกับคนที่อยากเข้าใจกลไกชั้นล่างของมือถืออย่างแท้จริง หรือแม้แต่อยากลองแฟลชรอมดู โปรดทราบว่าการทำงานเหล่านี้มีความเสี่ยงอยู่บ้าง โปรดดำเนินการด้วยความระมัดระวัง}

\booksection{หนึ่ง เข้าใจ Bootloader อย่างลึกซึ้ง}

Bootloader คือโปรแกรมแรกที่รันเมื่อมือถือเปิดเครื่อง หน้าที่ของมันคือเริ่มต้นฮาร์ดแวร์และโหลดเคอร์เนลของระบบปฏิบัติการ จะมอง Bootloader เป็น “ยามเฝ้าประตู” ของมือถือก็ได้

\booksubsection{หลักการทำงานของ BL Lock:}

\begin{itemize}
  \item \textbf{การตรวจสอบลายเซ็น}: BL Lock จะตรวจสอบลายเซ็นดิจิทัลของอิมเมจบูต มีแต่ระบบที่ผู้ผลิตลงลายเซ็นเท่านั้นถึงจะบูตได้
  \item \textbf{กันดาวน์เกรด}: กันผู้ใช้แฟลชระบบรุ่นเก่าลงไปเพื่อเจาะช่องโหว่ที่รู้จักแล้ว
  \item \textbf{ปกป้องข้อมูล}: แม้แฟลชระบบใหม่ลงไป ถ้าไม่มีรหัสผ่านบัญชีเดิมก็เข้าถึงข้อมูลผู้ใช้ไม่ได้
\end{itemize}

\booksubsection{การจัดเก็บและยืนยันระดับฮาร์ดแวร์}

สถานะการล็อก BL (Locked/Unlocked) และ Public Key ของผู้ผลิต มักถูกจัดเก็บไว้ในพื้นที่ฮาร์ดแวร์ที่ได้รับการปกป้อง (เช่น พาร์ติชัน devinfo ของ Qualcomm, eFuse หรือพาร์ติชัน RPMB) เพื่อป้องกันการถูกดัดแปลงด้วยซอฟต์แวร์ ในแต่ละครั้งที่บูต Bootloader จะอ่านสถานะจากพื้นที่ปลอดภัยนี้ และใช้เป็นเกณฑ์ในการตัดสินใจว่าจะดำเนินการตรวจสอบลายเซ็นแบบเข้มงวดหรือไม่

\booksubsection{การตรวจสอบลายเซ็นดิจิทัล}

เมื่อผู้ผลิตเผยแพร่อัปเดตระบบหรืออิมเมจ จะมีการเซ็นชื่อด้วย Private Key ของตนเอง เมื่ออุปกรณ์บูต Bootloader จะใช้ Public Key ในตัวเพื่อตรวจสอบอิมเมจระบบ (เช่น boot.img, recovery.img) วิธีการทำงานคือ:

\begin{enumerate}
  \item คำนวณค่าแฮชของอิมเมจระบบใหม่อีกครั้ง
  \item ใช้ Public Key ถอดรหัสลายเซ็นที่มากับอิมเมจ เพื่อให้ได้ค่าแฮชออกมาหนึ่งค่า
  \item เปรียบเทียบค่าแฮชทั้งสองค่า หากตรงกันแสดงว่าผ่านการตรวจสอบ ซึ่งจะรับประกันได้ว่าอิมเมจระบบไม่ถูกดัดแปลงโดยบุคคลที่สาม
\end{enumerate}

\booksubsection{ห่วงโซ่การบูตที่ปลอดภัย}

การล็อก BL คือวงแรกของห่วงโซ่การบูตที่ปลอดภัยทั้งหมด เมื่อผ่านการตรวจสอบแล้ว จึงจะโหลดและดำเนินการโปรแกรมบูตในขั้นตอนถัดไป (เช่น Kernel) เพื่อให้มั่นใจในความน่าเชื่อถือของสภาพแวดล้อมการบูตระบบทีละขั้นตอน

\booksubsection{ล็อก vs ปลดล็อก: ผลกระทบสองด้านของการล็อก BL}

เพื่อให้เข้าใจข้อดีและข้อเสียของการล็อก BL ได้ชัดเจนยิ่งขึ้น เราจะใช้ตารางเปรียบเทียบเพื่อแสดงความแตกต่างใน 5 มิติหลัก:

\begin{center}
\begin{tabularx}{\linewidth}{|X|X|X|}
\hline
\textbf{บทบาทหลัก} & \textbf{สถานะล็อก (ค่าเริ่มต้น)} & \textbf{สถานะปลดล็อก} \tabularnewline
\hline
ความปลอดภัยของระบบ & ป้องกันการแฟลชระบบที่ไม่เป็นทางการซึ่งไม่มีลายเซ็น ปกป้องความสมบูรณ์ของระบบและความปลอดภัยระดับล่าง\newline ความปลอดภัยสูง: โหลดเฉพาะอิมเมจที่ลงลายเซ็นอย่างเป็นทางการ ป้องกันการดัดแปลงเฟิร์มแวร์ที่เป็นอันตรายได้อย่างมีประสิทธิภาพ & ความปลอดภัยลดลง: อนุญาตให้แฟลช ROM ของบุคคลที่สามได้ กลไกการตรวจสอบความสมบูรณ์ของระบบล้มเหลว หากมัลแวร์ได้รับสิทธิ์ Root สามารถดัดแปลงระบบได้อย่างลึกซึ้ง \tabularnewline
ความเสถียรของระบบ & รับประกันว่าโทรศัพท์จะทำงานบนระบบทางการที่ผ่านการทดสอบและปรับแต่งโดยผู้ผลิตแล้ว เพื่อความเสถียรและลื่นไหล\newline เสถียร: ใช้เฟิร์มแวร์ทางการ ฟังก์ชันครบถ้วน อัปเดต OTA ที่เชื่อถือได้ & ความเสี่ยงเพิ่มขึ้น: ROM ของบุคคลที่สามอาจมีปัญหาความเข้ากันได้ อาจทำให้ระบบค้าง, อิฐ, หรือฟังก์ชันฮาร์ดแวร์ผิดปกติ (เช่น กล้อง, ชาร์จไว) การอัปเดต OTA อาจล้มเหลว \tabularnewline
สิทธิประโยชน์การรับประกัน & ผู้ผลิตล็อก BL เพื่อมอบระบบทางการที่เสถียน เพื่อให้บริการรับประกัน\newline รับประกันมีผล: ได้รับบริการรับประกันฟรีจากทางการ (กรณีไม่ใช่ความเสียหายจากผู้ใช้) & รับประกันสิ้นสุด: ผู้ผลิตส่วนใหญ่ถือว่าการปลดล็อก BL เป็นการสละสิทธิ์การรับประกันโดยสมัครใจ แม้จะล็อกกลับ ก็อาจถูกตรวจพบประวัติการปลดล็อกและปฏิเสธการรับประกัน \tabularnewline
ความปลอดภัยของข้อมูล & กระบวนการปลดล็อกมักจะบังคับล้างข้อมูล เพื่อป้องกันข้อมูลเก่าถูกนำไปใช้ในทางที่ผิดในสภาพแวดล้อมของระบบใหม่\newline ข้อมูลปลอดภัย: ในการใช้งานประจำวัน การล็อก BL จะจำกัดความสามารถของโปรแกรมที่เป็นอันตรายในการแก้ไขพาร์ติชันระบบ & ความเสี่ยงด้านข้อมูล: การปลดล็อกจะบังคับล้างข้อมูลผู้ใช้ทั้งหมด หากไม่เข้ารหัสอย่างถูกต้องหลังปลดล็อก การจัดเก็บข้อมูลอาจไม่ปลอดภัย หากสูญหายจะถูกแฟลชเพื่อดึงข้อมูลออกได้ง่ายขึ้น \tabularnewline
ข้อจำกัดของฟังก์ชัน & รับประกันว่าฟังก์ชันที่ต้องพึ่งพาสภาพแวดล้อมที่ปลอดภัย (เช่น การชำระเงิน, แอปธนาคาร) ทำงานได้ปกติ\newline ฟังก์ชันครบถ้วน: ฟังก์ชันทางการทั้งหมดสามารถใช้งานได้ปกติ & ฟังก์ชันจำกัด: แอปที่ละเอียดอ่อนด้านความปลอดภัยบางส่วน (เช่น Google Pay, แอปธนาคาร) อาจปฏิเสธการทำงานเนื่องจากตรวจพบสถานะการปลดล็อก BL บริการ DRM อาจถูกปรับลดระดับ ส่งผลกระทบต่อการเล่นวิดีโอความละเอียดสูง \tabularnewline
\hline
\end{tabularx}
\end{center}

\booksection{ขั้นตอนปลดล็อก BL โดยละเอียด (ยกตัวอย่าง Xiaomi)}

\begin{enumerate}
  \item \textbf{4. เริ่มต้นแฟลชรอม}รีเฟรชเครื่องใช่ชี้ให้โทรศัพท์มือถือติดตั้งใหม่ของดำเนินการระบบ ทั่วไปของรีเฟรชเครื่องประเภท:
  \item \textbf{รีเฟรชทางการROM}:ติดตั้งทางการระบบ, สำหรับกู้ภัยอิฐหรือการดาวน์เกรด
  \item \textbf{รีเฟรชบุคคลที่สามRecovery}:เหมือนTWRP, สำหรับสำรองข้อมูลกู้คืน, รีเฟรชเข้าบุคคลที่สามROM
  \item \textbf{รีเฟรชบุคคลที่สามROM}:เหมือนLineageOS, Pixel Experienceชั้น, ประสบการณ์ต่างกันของระบบ
\end{enumerate}

\booksection{วิธี Root หลายแบบ}

\begin{itemize}
  \item \textbf{รีเฟรชเครื่องมีความเสี่ยง, ดำเนินการต้องการระมัดระวัง อาจทำให้:}โทรศัพท์มือถือเปลี่ยนอิฐ(ไม่มีวิธีเปิดเครื่อง)
  \item \textbf{หมดประกัน}ข้อมูลทิ้งเสีย
  \item \textbf{จ่ายเงินการทำงานผิดปกติ}รีเฟรชเครื่องหน้าต้องทำดีเสร็จทั้งหมดสำรองข้อมูล, ระมัดระวังละเอียดอ่านอ่านสอนกระบวนการ, ยืนยันเครื่องรูปแบบและเวอร์ชันจับคู่
\end{itemize}

\booksection{เริ่มต้นแฟลชรอม}

XDA Developers

\begin{itemize}
  \item \textbf{: ชุมชนนักพัฒนา Android ที่ใหญ่ที่สุดในโลก มีบทช่วยสอนฟล้าเครื่องและแหล่ง ROM ให้เพียบ}Coolapk
  \item \textbf{: ชุมชนแกดเจ็ตดิจิทัลชื่อดังของจีน มีบทช่วยสอนภาษาจีนและกระดานอภิปรายมากมาย}ฟอรัมรายรุ่นมือถือ
  \item \textbf{: เช่น ชุมชน Xiaomi, ชุมชน OnePlus ฯลฯ มีทั้งบทช่วยสอนทางการและจากผู้ใช้}บทที่ 25: การจัดการมือถือมือสองและมรดกดิจิทัล
\end{itemize}

\booksection{Samsung Knox: ราคาที่ไม่สามารถย้อนกลับได้ของการปลดล็อก BL}

ต่างจากผู้ผลิต Android รายอื่น กลไกการปลดล็อก BL ของโทรศัพท์ Samsung มีดีไซน์ “ฟิวส์ขาด” ที่เป็นเอกลักษณ์:\textbf{สำหรับโทรศัพท์ Samsung การปลดล็อกจะทริกเกอร์การขาดของ Knox ทำให้สูญเสียฟังก์ชันความปลอดภัยบางส่วนอย่างถาวร (เช่น Samsung Pay); จะล้างข้อมูลที่ได้รับการปกป้องทั้งหมด; ส่งผลกระทบต่อการรับประกัน}โดยเฉพาะคือ:

\begin{itemize}
  \item \textbf{ฟิวส์นี้ใช้ได้ครั้งเดียว}: ภายในโทรศัพท์มีฟิวส์ eFuse อยู่หนึ่งตัว เมื่อตรวจพบว่า BL ถูกปลดล็อก มันจะขาดอย่างถาวรทันที ไม่สามารถกู้คืนได้
  \item \textbf{ฟังก์ชันที่ได้รับผลกระทบ}: Samsung Pay (บัตรโดยสาร/บัตรธนาคาร), โฟลเดอร์ปลอดภัย Knox, การจัดการ MDM ขององค์กร, แอปประกันสุขภาพบางส่วน จะใช้งานไม่ได้ทั้งหมด
  \item \textbf{บังคับล้างข้อมูล}: หลังจากทริกเกอร์การขาดของฟิวส์ โทรศัพท์จะบังคับรีเซ็ตเป็นค่าโรงงาน ข้อมูลในเครื่องทั้งหมดจะถูกลบ (ข้อมูลบนคลาวด์ไม่ได้รับผลกระทบ)
  \item \textbf{ผลกระทบต่อการรับประกัน}:ทางการระบุชัดเจนว่า หลังจาก Knox ถูกทริกเกอร์แล้ว จะไม่ได้รับบริการรับประกันอีกต่อไป
\end{itemize}

\begin{tcbwarn}{โปรดคิดให้ดีก่อนลงมือทำ}
Knox ของ Samsung ที่ถูกทริกเกอร์นั้น\textbf{ไม่สามารถย้อนกลับได้}เมื่อถูกทริกเกอร์แล้วจะไม่สามารถกลับสู่สถานะโรงงานเดิมได้ หากคุณแค่ต้องการใช้ฟังก์ชันที่ปกติต้องใช้ Root ขอแนะนำให้ตรวจสอบก่อนว่ามีทางเลือกอื่นที่ไม่ต้องใช้ Root หรือไม่

\end{tcbwarn}

\textbf{คำเตือนสำคัญ}\par

การแฟลชรอมมีความเสี่ยง ต้องดำเนินการด้วยความระมัดระวัง อาจทำให้:

\begin{itemize}
  \item : มือถือกลายเป็นอิฐ (เปิดเครื่องไม่ได้)
  \item เสียการรับประกัน
  \item ข้อมูลสูญหาย
  \item ฟังก์ชันชำระเงินทำงานผิดปกติ
\end{itemize}

: ก่อนแฟลชรอมต้องสำรองข้อมูลให้ครบถ้วน อ่านบทเรียนอย่างละเอียด ยืนยันว่ารุ่นเครื่องและเวอร์ชันตรงกัน

\begin{tcbinfo}{แหล่งข้อมูลที่แนะนำ}
\begin{itemize}
  \item \textbf{XDA Developers}: ชุมชนนักพัฒนา Android ที่ใหญ่ที่สุดในโลก มีบทเรียนแฟลชรอมและแหล่ง ROM ที่หลากหลาย
  \item \textbf{Coolapk}: ชุมชนดิจิทัลที่มีชื่อเสียงในจีน มีบทเรียนภาษาจีนและการอภิปรายมากมาย
  \item \textbf{ฟอรัมรายรุ่น}: เช่นชุมชน Xiaomi, ชุมชน OnePlus ฯลฯ มีทั้งบทเรียนจากทางการและภาคประชาชน
\end{itemize}

\end{tcbinfo}

\booksection{คำถามที่พบบ่อย (Q\& A)}

\textbf{Q1: หลังจากปลดล็อก Bootloader แล้วสามารถกลับไปล็อกใหม่ได้ไหม?}\par

\textbf{A:}สามารถล็อกใหม่ได้ โดยปกติจะต้อง Flash ชุดเฟิร์มแวร์อย่างเป็นทางการ และใช้สคริปต์ Flash เฉพาะ (เช่น ของ Xiaomi\texttt{flash\_all\_lock.\allowbreak bat}) เพื่อล็อกใหม่ได้ แต่โปรดทราบว่า หลังจากล็อกใหม่แล้ว การทริกเกอร์ความปลอดภัยบางอย่างที่เกิดจากการปลดล็อก (เช่น Knox ของ Samsung) อาจไม่สามารถย้อนกลับได้

\textbf{Q2: หลังจากปลดล็อกแล้วจะ Root ได้อย่างไร?}\par

\textbf{A:}การปลดล็อก Bootloader เป็นเพียงก้าวแรก การ Root มักจะต้องทำดังนี้:

\begin{enumerate}
  \item Flash Recovery ที่กำหนดเอง (เช่น TWRP)
  \item Flash Magisk หรือตัวจัดการ Root อื่นๆ ผ่าน Recovery เพื่อแก้ไข boot image จึงจะได้สิทธิ์ Root
\end{enumerate}
